%==============================================================================
% Appendix C: Full Shear Tensor Basis and Embedding of (eps, s, phi)
%==============================================================================
\section{Full Shear Tensor Basis and Embedding of $(\varepsilon,s,\phi)$}
\label{app:C}

This appendix summarises the final basis of the spin-2 sector and its relation
to the raw geometric variables $(\varepsilon,s,\phi)$.

\subsection{Traceless Symmetric Basis}

The full shear tensor is expanded as
\begin{equation}
  h_{ij} = \sum_{A=1}^5 q_A\,(T_A)_{ij}.
\end{equation}
The explicit matrix representations of $T_1,\ldots,T_5$ are given in
\S\ref{sec:setup:spin2}.  The following properties hold exactly:
\begin{equation}
  \tr(T_A T_B) = \delta_{AB},
  \qquad
  \tr(h) = 0,
  \qquad
  \tr(h^2) = \sum_{A=1}^5 q_A^2.
\end{equation}

\subsection{Embedding of the Old Variables}

The raw geometric variables $(\varepsilon,s,\phi)$ provide a description of a
particular three-dimensional slice within the spin-2 quintet.  In this paper we
do not place these slice variables in the primary role; instead we adopt the
final basis $q_A$ as the principal variables.  Accordingly:
\begin{itemize}
  \item $\varepsilon$ is primarily used for axisymmetric deformations in the
        $m=0$ singlet direction;
  \item $s$ and $\phi$ are auxiliary parameters for tracking the negative-side
        condensed phase.
\end{itemize}
All representation-theoretic statements are always rewritten in $q_A$.

\subsection{Invariants in the Final Basis}

It is confirmed exactly that the low-order invariants in the final basis are
\begin{equation}
  I_2 = \tr(h^2) = \sum_{A=1}^5 q_A^2,
  \qquad
  I_3 = h_{ij}\omega_i\omega_j,
  \qquad
  I_4 = \tr(h^3).
\end{equation}

\subsection{Constant-Shape Homogeneous Consistency}

The homogeneous consistency for the constant frame deformation
$e^i = r\, M^i{}_j\,\sigma^j$ with the structure-constant transformation
\begin{equation}
  C'^i{}_{jk} = M^i{}_a\,(M^{-1})^b{}_j\,(M^{-1})^c{}_k\, C^a{}_{bc}
\end{equation}
was verified for two representative all-five-nonzero full-shear samples on both
$\Sthree$ and $\Nilthree$:
\begin{itemize}
  \item antisymmetry residual $= 0$,
  \item Jacobi residual $= 0$,
\end{itemize}
both holding exactly.  The final spatial ansatz is therefore self-consistent as a
homogeneous minisuperspace frame.

\subsection{Role of This Appendix}

The purpose of this appendix is to clarify the change of viewpoint from raw
geometric variables to the final basis.  All final statements about spin-2 in
this paper are written in $q_A$ and its irreducible decomposition; $(\varepsilon,
s,\phi)$ are used only as auxiliary parameters for the phase diagram.

\subsection{Jacobian: Old Variables to Final Basis}

We give the correspondence between the raw geometric variables $(\varepsilon,s)$
and the final basis coordinates $(q_1,q_2)$ explicitly.

For $\varepsilon$---which embeds a $T_1$-direction deformation exponentially:
\begin{equation}
  \varepsilon = \exp\!\left(\frac{\sqrt{6}}{2}\,q_1\right) - 1,
  \qquad
  q_1 = \frac{2}{\sqrt{6}}\ln(1+\varepsilon),
  \qquad
  \frac{\partial q_1}{\partial\varepsilon}\bigg|_{\varepsilon=0}
  = \frac{2}{\sqrt{6}} = \sqrt{\frac{2}{3}}.
\end{equation}

For $s$ (deformation in the $T_2$ direction):
\begin{equation}
  s = \exp\!\left(\frac{q_2}{\sqrt{2}}\right) - 1,
  \qquad
  q_2 = \sqrt{2}\ln(1+s),
  \qquad
  \frac{\partial q_2}{\partial s}\bigg|_{s=0} = \sqrt{2}.
\end{equation}

The angle $\phi$ appearing in the $s\neq 0$ phase corresponds to an $SO(2)$
rotation in the $(q_2,q_3)$ plane:
\begin{equation}
  q_2 = s_q\cos\phi,
  \qquad
  q_3 = s_q\sin\phi,
  \qquad
  s_q = \sqrt{q_2^2+q_3^2}.
\end{equation}
The analogous embedding for the $T_4,T_5$ directions (off-diagonal components)
is structurally identical by symmetry.

The combined Jacobian diagonal at the isotropic point is
\begin{equation}
  J = \mathrm{diag}\!\left(\sqrt{\tfrac{2}{3}},\,\sqrt{2}\right),
\end{equation}
and the Hessian basis transformation is
\begin{equation}
  H_q = J^{-T} H_{\varepsilon s}\, J^{-1}.
\end{equation}
At the isotropic point, $H_{\varepsilon s} = \mu_q^2\,I_2$ and
$J^{-T}J^{-1}\propto I_2$, so Schur degeneracy $H_q = \mu_q^2\,\delta_{AB}$
is preserved under the change of basis.
